\section{Introduction}
\label{sec:intro}

A landmine detector that is evaluated with mean average precision is being asked
the wrong question. Mean average precision integrates performance over every
confidence threshold and every recall level, then reduces the result to one
number. A demining team does not operate a curve; it operates one threshold, on
one day, with one tolerance for false alarms. Two systems can share a curve and
behave completely differently at the threshold that will actually be used. This
paper is an extended argument, with measurements, that the operating point---and
not the summary curve---is the object that should be reported, and that the
usually unglamorous work of auditing the split is what makes any such comparison
believable.

We study single-class detection of \emph{surface} landmine indicators from
downward-looking aerial imagery. The scope limitation is physical and we state it
first: a visible-light or colour camera cannot see a buried object. What it can
see is a mine lying on the surface, a partially exposed casing, or a surface
marker. Claims in this paper are therefore claims about visible surface objects,
not about buried ordnance, and nothing here is a field-validated clearance
system. The task is nonetheless a useful stress test for detection methodology,
because the class is rare, the imagery is heterogeneous, and the published data
reuse is heavy.

\paragraph{What we found.} Three results organise the paper.

\begin{enumerate}[leftmargin=1.2em,itemsep=2pt]
\item \textbf{The split leaks, and the leak is measurable.} The combined
landmine corpus contains $76$ filename groups (Roboflow base names) that span
train and evaluation splits, and $63.8\%$ of validation and test images have a
perceptual-hash neighbour in training. This is the signature of augmentation
applied before splitting: the same scene, re-exported, lands on both sides of the
boundary. After rebuilding the split group-wise and filtering twice, $30$
evaluation images are removed as true leaks. A detector trained on the leaky
split scores $0.994$ mAP$_{50}$ on the full clean test set, against $0.974$ for a
clean-trained detector. On the $49$ test images that are contamination-free for
both, the two are tied at $0.995$ and $0.995$. The leak inflated the number, but
by less than the folklore predicts---a negative result, and one that would have
been invisible without the contamination-free subset.

\item \textbf{mAP$_{50}$ does not settle the deployment question.} A YOLOv8s
detector over three seeds reaches $0.977 \pm 0.003$ test mAP$_{50}$; the one
available YOLOv8n seed reaches $0.979$. Although the n point is slightly higher
on test, validation mAP favors s ($0.954$ vs. $0.932$), and n has no seed
interval. At the recall-first operating point chosen on validation, the observed
n checkpoint produces $6.50$ false positives per image against $1.51$ for s.
We present this as a case study in threshold transfer and operating-point
reporting, not as a reliable model-ranking reversal.

\item \textbf{The two-stage gain is detector-dependent, and we can say why.} A
MobileNetV3-Large classifier over detector crops improves precision over a
recall-matched single-stage baseline at two of three detector seeds
($\Delta P = +0.128$ and $+0.144$, paired bootstrap intervals excluding zero).
At the third seed it is \emph{worse} ($\Delta P = -0.037$, interval
$[-0.077, +0.013]$). Under a max-$F_2$ rule, a single confidence threshold tuned
on validation beats the two-stage pipeline outright ($P = 0.976$ vs.\ $0.945$).
The pattern is consistent with the classifier helping when the validation-selected
detector threshold transfers poorly to test, but three seeds do not identify that
mechanism. Where the baseline threshold transfers well, the second stage has
nothing to add. Reporting a three-seed mean of $0.879 \pm 0.018$ would have hidden
this entirely; we report the per-seed structure instead.
\end{enumerate}

\paragraph{Why this is worth publishing.} None of the three findings required a
new architecture, a large GPU, or a new dataset. They required (i) reconstructing
the split from scratch, (ii) refusing to average over seeds when the seeds
disagree structurally, and (iii) measuring at the operating point rather than at
the summary. The total compute was one 8\,GB laptop GPU and roughly $30$ hours.
We think this is a realistic and under-reported budget for applied detection
work, and we document what it forbade (Section~\ref{sec:limitations}) rather than
implying broader coverage.

\paragraph{Contributions.}
\begin{itemize}[leftmargin=1.2em,itemsep=2pt]
\item A reproducible leakage audit of a publicly derived landmine corpus: an
operational leak definition, a two-method filter with a manually validated MSE
cut, a frozen manifest, and an accounting of all $30$ removals
(Section~\ref{sec:audit}).
\item A direct leaky-versus-clean comparison on a contamination-free subset,
including the unfavourable result that the mAP difference on a small
contamination-free subset is small, while threshold-matched operating-point
effects remain unmeasured
(Section~\ref{sec:leakvalue}).
\item An operating-point evaluation protocol---thresholds chosen on validation
only, applied once to test, with paired image bootstrap and source-cluster
intervals---and a case study showing why a single mAP-based ranking can be
insufficient
(Section~\ref{sec:results}).
\item A per-seed analysis of a two-stage detector-plus-classifier pipeline that
localises a failure rather than averaging it away, together with a max-$F_2$
analysis showing when the second stage is not worth its cost
(Section~\ref{sec:twostage}).
\item A Grad-CAM audit used as a sanity check on evidence, with an explicit
separation between attribution and correctness, showing that the image-level
classifier's accuracy is not accompanied by object-level evidence
(Section~\ref{sec:gradcam}).
\end{itemize}

Section~\ref{sec:related} positions the work; Section~\ref{sec:data} describes the
corpus and its provenance; Section~\ref{sec:audit} the audit;
Sections~\ref{sec:method}--\ref{sec:protocol} the models and protocol;
Section~\ref{sec:results} the results; Sections~\ref{sec:discussion}--\ref{sec:ethics}
discussion, limitations and ethics. Code, weights, tables and the frozen split
manifest accompany the paper (Section~\ref{sec:availability}).
